% Dokumentanklasse: a4paper, 14pt
% Beschreibung:     Dokumentenformat
% Option:           extraarticle - ?
\documentclass[12pt]{article}

% Paket:            setspace
% Beschreibung:     Setz über die optionen den Zeilenabstand
% Optionen:         Möglicher Zeilenabstand
%                   singlespacing = 1,0
%                   onehalfspacing = 1,5
%                   doublespacing = 2,0
% Restriktion:      Muss von Paket hyperref geladen werden. Ansonsten funktioniert das Paket nicht.
\usepackage[onehalfspacing]{setspace}

% Paket:            appendix
% Beschreibung:     Das Paket dient dazu, ausschließlich das Thema einer Überschrift in das Inhaltsverzeichnis zu überführen
% Option:           appendix - Überführt die Überschriften des Anhangs richtig ins das Inhaltsverzeichnis
\usepackage[titletoc]{appendix}

% Paket:            ansmath
% Beschreibung:     Zum darstellen von mathematischen Formeln
\usepackage{amsmath}

% Paket:            bebel
% Beschreibung:     Stellt das Literatur-, Abbilungs-, Tabellenverzeichnis auf eine beliebige Sprache
% Option:           ngerman
\usepackage[ngerman]{babel}

% Paket:            biblatex
% Beschreibung:     Ermöglicht es, ein Literaturverzeichnis zu führen
% Option:
\usepackage[
  style=authoryear-icomp,    % Zitierstil
  isbn=false,                % ISBN nicht anzeigen, gleiches geht mit nahezu allen anderen Feldern
  pagetracker=true,          % ebd. bei wiederholten Angaben (false=ausgeschaltet, page=Seite, spread=Doppelseite, true=automatisch)
  maxbibnames=50,            % maximale Namen, die im Literaturverzeichnis angezeigt werden (ich wollte alle)
  maxcitenames=3,            % maximale Namen, die im Text angezeigt werden, ab 4 wird u.a. nach den ersten Autor angezeigt
  autocite=inline,           % regelt Aussehen für \autocite (inline=\parancite)
  block=space,               % kleiner horizontaler Platz zwischen den Feldern
  backref=true,              % Seiten anzeigen, auf denen die Referenz vorkommt
  backrefstyle=three+,       % fasst Seiten zusammen, z.B. S. 2f, 6ff, 7-10
  date=short,                % Datumsformatbackend=biber
]{biblatex}
\setlength{\bibitemsep}{1em}     % Abstand zwischen den Literaturangaben
\setlength{\bibhang}{2em}        % Einzug nach jeweils erster Zeile
\addbibresource{referenzen//bibliothek.bib}

% Paket:            caption
% Beschreibung:     Bietet eine große Auswahl an Gestaltungsmöglichkeiten bezüglich der Beschriftung von Bildern und Tabellen.
\usepackage{caption}

% Paket:            subcaption
% Beschreibung:     Bietet eine große Auswahl an Gestaltungsmöglichkeiten bezüglich der Beschriftung von Bildern und Tabellen
%                   die neben oder unter untereinander dargestellt werden sollen.
\usepackage{subcaption}

% format=NAME           - Einstellung des grundsätzlichen Formats von Gleitobjektbeschriftungen
% indention=EINZUG      - Einstellung des Einzugs der Beschriftung ab der zweiten Zeile
% labelformat=NAME      - Einstellung der Zusammensetzung des Bezeichners (z. B. Label und Nummer) der Beschriftung
% labelsep=NAME         - Einstellung des Trenners nach dem Bezeichner
% justification=NAME    - Einstellung er Ausrichtung des Textes Beschriftung
% singlelinecheck=BOOL  - Sonderbehandlung für einzeilige Beschriftungen ein- oder ausschalten
% font=NAME             - Einstellung der Schrift der gesamten Beschriftung
% labelfont=NAME        - Einstellung der Schrift des Bezeichners
% textfont=NAME         - Einstellung der Schrift des Textes der Beschriftung
% margin=BREITE         - Einstellung der Breite des Randes der Beschriftung
% width=BREITE          - Einstellung der Breite der Beschriftung
% parskip=ABSTAND       - Einstellung des Absatzabstandes der Beschriftung
% hangindent=EINZUG     - Einstellung des Einzugs abgesehen von ersten Absatzzeilen
% style=NAME            - Auswahl einer vordefinierten Gesamteinstellung
% aboveskip=ABSTAND     - Einstellung des Abstandes vor einer Beschriftung
% belowskip=ABSTAND     - Einstellung des Abstandes nach einer Beschriftung
% position=WAHL         - Einstellung ob die Beschriftung als Über- oder Unterschrift formatiert werden soll
% tableposition=WAHL    - Einstellung ob nur bei Tabellen die Beschriftung als Über- oder Unterschrift formatiert werden soll
\captionsetup[figure]{position=bottom}
\captionsetup[table]{position=bottom}

% Paket:            courier
% Beschreibung:     Lädt das Paket courier für Schriftarten mit fester Breite.
% Befehle:          \ttfamily     Aktiviert Courier füt Tabellen bzw. generelle begin-Blöcke
\usepackage{courier}

% Paket:            enumitem
% Beschreibung:     Zeilenabstände bei Aufzählungen definieren
% Option:
\usepackage{enumitem}

% Paket:            eurosym
% Beschreibung:     Bildet das Euro-Zeichen in unterschiedlichen Varianten ab
% Option:
\usepackage{eurosym}

% Paket:            fancyhdr
% Beschreibung:     Ermöglich ein generelles Seitenlayout ein zu stellen mit Kopf und Fußzeile.
\usepackage{fancyhdr}

% Paket:            float
% Beschreibung:     Zum Ausrichten von Tabellen und Spalten bzw. deren Zentrierung
% Option:
% Restriktion:      Muss von Paket hyperref geladen werden. Ansonsten funktioniert das Paket nicht.
\usepackage{float}

% Packet:           framemethod
% Beschreibung:
% Option:
\usepackage[framemethod=tikz]{mdframed}
\mdtheorem[
  linecolor=red,
  frametitlefont=\sffamily\bfseries\color{white},
  frametitlebackgroundcolor=red,
]{warn-popup}{Warnung}[subsection]

\mdtheorem[
  linecolor=orange,
  frametitlefont=\sffamily\bfseries\color{white},
  frametitlebackgroundcolor=orange,
]{info-popup}{Information}[subsection]

\mdtheorem[
  linecolor=green,
  frametitlefont=\sffamily\bfseries\color{white},
  frametitlebackgroundcolor=green,
]{example-popup}{Beispiel}

% Paket:            geometry
% Beschreibung:     A4 Seitenabstände
% Option:
\usepackage{geometry}
\geometry{
%  a4paper,           % Papierformat (wird auch über die Dokumentenklasse definiert)
  top=2.7cm,          % Abstand Kopfseite   (Zwischen Kopfseite und Inhalt)
  bottom=2cm,         % Abstand Fußseite    (Zwischen Fußseite und Inhalt)
  left=2.5cm,         % Abstand Linkeseite  (Zwischen Linkerseite und Inhalt)
  right=2cm,          % Abstand Rechteseite (Zwischen Recherseite und Inhalt)
%  width=              %                    textwidth (+marginsep +marinparwidth)
%  textwidth=15cm,     % Textbreite
%  marginparsep=1cm,   % Randnotiztrennng
%  marginparwidth=10cm,% Randnotizbreite
%  height=             %                    textheight (+headheight +headsep + footskip)
%  textheight=        % Texthöhe
  headheight=1cm,     % Kopfhöhe
  headsep=0.5cm,      % Kopftrennung        (Größe zwischen Kopfzeile und Inhalt)
  footskip=1cm,       % Fußzeilenhöhe
}

% Paket:            glossaries
% Beschreibung:
% Option:
\usepackage[toc,acronyms]{glossaries}
\makeglossaries
% Akronyme
\newacronym{acr:sql}{SQL}{Structured-Query-Language}
\newacronym{acr:ddl}{DDL}{Data-Definition-Language}

% Glossar
\newglossaryentry{gls:minted}{
  name={Minted},
  description={Minted is a package that allows formatting and highlighting source code in \latex.}
}



% Paket:            graphicx
% Beschreibung:     Einbinden von Bildern
% Option:
\usepackage{graphicx}

% Packet:           Hyperref
% Beschreibung:     Importiert hyperref um Querverweise mittels \hyperref zu erzeugen.
\usepackage{hyperref}
\hypersetup{
  pdftitle={PDF-Title},
  pdfauthor={Markus Pesch},
  pdfsubject={PDF-Subject},
  pdfkeywords={},
  pdfcreator={pdflatex},
  pdfproducer={LaTeX with hyperref}
}

% Paket:            utf8
% Beschreibung:     Stellt Umlaute richtig dar
% Option:           inputenc - Erlaubt die Darstellung der gleichen Zeichen (Character) wie sie in stdin überliefert werden
\usepackage[utf8]{inputenc}

% Paket:            makeindex
% Beschreibung:     Ermöglicht das Indexieren von Wörter und den Befehl \printindex um den Index auszugeben
\usepackage{makeidx}
\makeindex

% Packet:           Minted
% Beschreibung:     Dient zum highlining von Quellcode wie beispielsweise Java, Bash oder Python.
% Option/en:
%   autogobble:       Leerzeichen zwischen linken Rand und Sourcecode einrücken bzw. weg schneiden.
%   breaklines:       Automatische Zeilenumbrüche
%   cache:            de- oder aktiviert den cache um Sourcecode zwischen zu speichern und so das PDF schneller zu erzeugen
%   cachedir:         Definiert den Pfad zum cache, an dem minted seine Daten zwischen speichern kann
%   fontfamily:       Die Schriftart die benutzt werden soll. tt, courier und helvetica sind vordefiniert.
%   fontsize:         Die Schriftgröße die benutzt werden soll. Beispielsweise fontsize=\footnotesize
%   linenos:          Zeilennummern
%   keywordcase:      Änderung der Buchstaben. Takes lower, upper, or capitalize.
%   showspaces:       Blendet Leerzeichen ein
\usepackage[cache=true, cachedir=/tmp/minted]{minted}

% usemintedstyle
% Gebe 'pygmentize -L styles' im Terminal ein um alle verfügbaren styles anzuzeigen.
\usemintedstyle{tango}

% newminted
% Definiere neue aliase um einmalig ein highlighting pro Sprache zu deklarieren
% \newminted{<makroname>}{optionen} ist verfügbar unter "<makroname>code"
\newminted{awk}{autogobble=true, breaklines=true, linenos=true}
\newminted{json}{autogobble=true, breaklines=true, linenos=true}
\newminted{r}{autogobble=true, breaklines=true, linenos=true}
\newminted{sql}{autogobble=true, breaklines=true, linenos=true, keywordcase=upper}
\newminted{xml}{autogobble=true, breaklines=true, linenos=true}

% newmintedfile
% Definiere neue Makros um automatisch Sourcecode aus Dateien zu highlighten.
% \makroname{Dateipfad}
\newmintedfile[inputawk]{awk}{autogobble=true, breaklines=true, linenos=true}
\newmintedfile[inputjson]{json}{autogobble=true, breaklines=true, linenos=true}
\newmintedfile[inputr]{r}{autogobble=true, breaklines=true, linenos=true}
\newmintedfile[inputsql]{sql}{autogobble=true, breaklines=true, linenos=true, keywordcase=upper}
\newmintedfile{inputxml}{autogobble=true, breaklines=true, linenos=true}

% newmintinline
% Definiere neues Makro um Sourcecoude einzeiler zu highlighten
% \begin{awkcode} \end{awkcode}
\newmintinline{awk}{autogobble=true, breaklines=true, linenos=true}
\newmintinline{json}{autogobble=true, breaklines=true, linenos=true}
\newmintinline{r}{autogobble=true, breaklines=true, linenos=true}
\newmintinline{sql}{autogobble=true, breaklines=true, linenos=true, keywordcase=upper}
\newmintinline{xml}{autogobble=true, breaklines=true, linenos=true}

% Paket:            multirow
% Beschreibung:     Zum kombinieren mehrerer Zellen einer Tabelle
% Option:
\usepackage{multirow}

% Paket:            natbib
% Beschreibung:     Für Zitate
% Option:           round - ?
%\usepackage[round]{natbib}

% Paket:            pdflscape
% Beschreibung:     Ermöglicht Seiten horizontal darzustellen
% Option:           \begin{landscape} \end{landscape}
\usepackage{pdflscape}

% Paket:            showframe
% Beschreibung:     Blendet alle Frames (Textkörper, Fußzeile Kopzeile, Seitenrand) ein
% Option:
% \usepackage{showframe}

% Packet:           tabularx
% Beschreibung:     Werden Tabellen mit diesem Paket erstellt, ist es möglich Zeilenumbrüche in einer Zelle zu erzeugen
\usepackage{tabularx}

% Paket:            textpos
% Beschreibung:     Zum erstellen von Textboxen
% \begin{textblock*}{\textwidth}(0cm,0.5cm)
\usepackage{textpos}

% Paket:            verbatim
% Beschreibung:     Bildet einen Quelltext exakt ab.
% Options:
\usepackage{verbatim}

% Paket:            warapfig
% Beschreibung:     Ermöglich das floaten von Texten neben Bildern
% Options:
% \begin{wrapfigure}{R}{0.30\textwidth}
%   \includegraphics[width=0.30\textwidth]{img/middleware.png}
%   \caption{Middleware}
% \end{wrapfigure}
\usepackage{wrapfig}

% Packet:           xcolor
% Beschreibung:     Define own color schemas
% Option:
\usepackage{xcolor}

% Definiere Farben
\definecolor{blue}{HTML}{2E579F}
\definecolor{green}{HTML}{1B8347}
\definecolor{orange}{HTML}{FF7F00}
\definecolor{red}{HTML}{D31C27}
\definecolor{white}{HTML}{FFFFFF}

% Packet:           csquotes
% Beschreibung:     Muss nach babel, polyglossia, biblatex und inputec geladen werden
% Option:
\usepackage{csquotes}

% Einstellungen überschreiben
\newcommand{\myparagraph}[1]{\paragraph{#1}\mbox{}\\}

% Start des Dokuments
\begin{document}

  % Set Globally Table-Margin
  \def\arraystretch{1.2}

  % Fetch Commit ID and Date
  \immediate\write18{./git-info.sh commit > git-id.tmp}
  \immediate\write18{./git-info.sh date > git-date.tmp}
  \immediate\write18{./git-info.sh url > git-url.tmp}

  % Importiere weitere .tex Dokumente
  \begin{titlepage}
  \begin{center}
    \begin{large}
      Tutorium SS18
    \end{large}
    
    \begin{huge}
      \begin{singlespace}
            \textbf{Datawarehouse}
      \end{singlespace}
    \end{huge}

    \vspace{0.5cm}

    \begin{figure}[h]
      \centering
      \includegraphics[width=0.85\textwidth]{img//logo.png}
      \label{img:fh-trier-logo}
    \end{figure}

    \vspace{2cm}
    
    Latex Sourcecode on \input{./git-url.tmp} \\
    based on git commit \input{git-id.tmp} from \input{git-date.tmp}
    
    
  \end{center}
\end{titlepage}

  \pagebreak

  % Pagestyle
  % Setze das Seitenlayout auf fancyhdr um Fuß- und Kopfzeilen zu setzen
  \pagestyle{fancy}

  % Löscht alle Kopf- und Fußzeilen des pagestyles fancyhdr
  \fancyhf{}

  % Fuß- und Kopfzeile des Paketes fancyhdr
  % [L] - Linkeseite      [O] - Ungerade Seitenzahlen         [LE,LO] - Linkeseite, Gerade- und Ungerade Seitenanzahlen
  % [C] - Mitte           [E] - Gerade Seitenanzahlen         [CE]    - Seitenmitte, nur gerade Seitenanzahlen
  % [R] - Rechteseite                                         [RO]    - Rechteseite, nur ungerade Seitenanzahlen
  % \fancyhead    Kopfzeile
  % \fancyfoot    Fußzeile
  \fancyhead[L]{\rightmark}
  \fancyhead[R]{\includegraphics[width=4cm]{img/logo.png}}
  \fancyfoot[L]{Seminar Data Science SS18 - Git}
  \fancyfoot[C]{}

  % Pixelstärke der Kopf- und Fußzeilenlinie
  \renewcommand{\headrulewidth}{1pt}
  \renewcommand{\footrulewidth}{1pt}

  % Agenda
  \tableofcontents
  \pagebreak

  % Setze die Seitenbeginn zurück
  \setcounter{page}{1}
  \fancyfoot[R]{Seite \thepage}

  % Importiere weitere .tex Dokumente
  \section{Einführung}
\label{sec:einfuehrung}
Dieses Kapitel bietet eine Einführung in die Grundbegriffe von Git und deren generelle Grundkonfiguration.

\subsection{Grundbegriffe}
\label{sec:einfuehrung.grundbegriffe}
Um dem Umgang mit Git zu erläutern ist es notwendig, einige Fachbegriffe zu verstehen.

\myparagraph{Versionskontrollsystem}
Ein Versionskontrollsystem dient zur Verwaltung und Versionierung von Software. Bekannte Projekte die unter dem Versionskontrollsystem Git stehen sind beispielsweise der Linux-Kernel und Git selbst. Beide Projekte stehen unter der Lizenz GPL. Der Quellcode kann mittels git bezogen werden. Dazu später mehr.

Bei Versionskontrollsystemen wird unterschieden zwischen zentralen und dezentralen Systemen. Git ist ein dezentrales System. Der Entwickler ist nicht abhängig von dem Server auf dem die Versionshierarchie gespeichert ist. Der Entwickler ist in der Lage mit Git die komplette Versionshierarchie auf sein eigenes System zu klonen und auf dieser Basis zu arbeiten.

\myparagraph{Repository}
Git  verwaltet von jeder Datei unterschiedliche Zustände bzw. Versionen. Diese werden in dem Repository gespeichert. Mithilfe des Repositories ist es möglich, auf jeden Zustand einer Datei zurück zu springen.

\myparagraph{Working Tree}
Alle Modifikationen an Dateien bzw. dem Quellcode werden im Working Tree vorgenommen. Andere Bezeichnungen hierfür sind auch Working Directory oder Workspace. Wobei bei vielen IDEs im Heimatverzeichnis das Verzeichnis workspace erstellt wird um Projekte zu speichern und unter Versionskontrolle zu setzen.

\myparagraph{Commit}
Ein Commit nimmt jede Veränderung, Modifikation oder Erstellung an Dateien im Working Tree auf. Zusätzlich speichert der Commit noch weitere Informationen wie den Benutzernamen und dessen E-Mail mit Datum und optional gpg-Signatur ab.

\myparagraph{HEAD}
HEAD bildet eine Referenz ab, in welchem Zustand der Entwickler sein Working Tree vor findet.

\myparagraph{Branch}
Jede Software kann in ihrem Verlauf mehrere Entwicklungsverzweigungen haben um beispielsweise in jeder Verzweigung einzelne Teilaufgaben wie Patches oder neuen Features von Software zu entwickeln. Diese Verzweigungen werden Branches genannt. Jede Verzweigung kann später durch einen Merge oder Rebase wieder zurück geführt werden.

\subsection{Grundeinrichtung}
\label{sec:einfuehrung.git}
Wie bereits in in Kapitel \ref{sec:einfuehrung.grundbegriffe} erklärt, speichert Git bei jedem Commit den Benutzernamen und E-Mail ab um nachzuvollziehen von wem die Änderung des Quellcodes stammt. Bevor nun dieses Dokument, das als \LaTeX{} Projekt angelegt wurde, durch Git verwaltet wird, ist es Sinnvoll den Benutzernamen und die E-Mail Adresse anzupassen \footcite{git-1.6-your-identity}.

\begin{minted}[framesep=2mm, fontsize=\small]{bash}
$ git config --global user.name "Hugo McKinnock"
$ git config --global user.email "hugo.mckinnock@example.com"
\end{minted}

\subsection{Signieren von Commits}
\label{sec:einfuehrung.sign-commits}
Da jede Person in der Lage ist den Quellcode boshaft zu ändern unter beliebigen Namen, bietet Git die Möglichkeit an jeden Commit mit einem GnuPG Schlüssel zu signieren.

\begin{info-popup}
  Manche Projekte setzen signierte Commits voraus, wie beispielsweise das Projekt um den Linux Kernel. Nicht signierte Commits werden durch die Hauptverantwortlichen nicht übernommen. Dazu später mehr.
\end{info-popup}

Ist GnuPG auf dem eigenen System nicht vorhanden, kann dies nachträglich installiert werden. Debian basierende Derivate  bieten GnuPG durch die offiziellen Paketquellen an. Dazu reicht die Installation über den Terminalbefehl \textit{sudo apt-get update; sudo apt-get install gpg2 --yes}.

Wird ein Windows-System genutzt, müssen die GnuPG Binaries manuell heruntergeladen und installiert werden. Die Binaries können unter \href{https://www.gpg4win.de}{gpg4win.de} bezogen werden.

Der GnuPG Schlüssel wird anschließend über den Aufruf von gpg im Terminal (Linux Derivate) oder der Git-Bash (Windows) erzeugt. Im Beispiel wird ein RSA 2048 Bit Schlüssel erzeugt.

\begin{verbatim}
$ gpg --gen-key
gpg (GnuPG) 2.1.11; Copyright (C) 2016 Free Software Foundation, Inc.
This is free software: you are free to change and redistribute it.
There is NO WARRANTY, to the extent permitted by law.

Hinweis: "gpg --full-gen-key" ruft den erweiterten Dialog auf.

GnuPG erstellt eine User-ID, um Ihren Schlüssel identifizierbar zu machen.

Ihr Name: Hugo McKinnock
Email-Adresse: hugo.mckinnock@example.com
Sie haben diese User-ID gewählt:
"Hugo McKinnock <hugo.mckinnock@example.com>"

Ändern: (N)ame, (E)-Mail oder (F)ertig/(A)bbrechen? F
gpg: Schlüssel F3FFC9EC ist als ultimativ vertrauenswürdig gekennzeichnet
\end{verbatim}

Wurde der GnuPG Schlüssel erzeugt, kann dieser nun hinterlegt werden um Commits im Hintergrund automatisch zu signieren. Dazu muss zuerst die Schlüssel-ID ausfindig gemacht werden und anschließend Git mitgeteilt werden.

\begin{verbatim}
$ gpg --list-secret-keys --keyid-format 0xSHORT
/home/hugo/.config/gnupg/pubring.gpg
--------------------------------------
sec   rsa2048/F3FFC9EC 2018-03-15 [SC]
uid      [  ultimativ] Hugo McKinnock <hugo.mckinnock@example.com>
ssb   rsa2048/4A1B6E12 2018-03-15 [E]

$ git config --global user.signingkey "0x85ED78DE"
$ git config --global commit.gpgSign true
\end{verbatim}

\begin{info-popup}
  Sollte das Betriebssystem gpg nicht finden, wird die Binary nicht durch die Umgebungsvariable \textit{PATH} aufgelöst. Entweder wird der Pfad zur gpg Binary der Umgebungsvariable \textit{PATH} hinzugefügt oder der absolute Pfad mit Angabe der gpg Binary Git durch folgenden Befehl mitgeteilt - \textit{git config --global gpg.program $ < $Pfad$ > $}
\end{info-popup}

  % Aufzählung aller Bilder
  % \listoffigures
  % \newpage

  % Aufzählung aller Tabellen
  % \listoftables
  % \newpage

  % Glossary
  \printglossaries

\end{document}

